% Allocation response proofs and numerical comparisons; CPU-verified.
\section{Frequency Allocation and Distance-Dependent Content Responses}
\label{sec:allocation-response}

\subsection{Positional overlap preserves content coordinates}
\label{sec:content-coordinate-retention}
For a block $\mathcal B$ of $r$ rotary pairs, write
$\mathcal R_{\mathcal B}(d)=\operatorname{diag}_{k\in\mathcal B}R(\omega_kd)$.
Each block is orthogonal, so $\mathcal R_{\mathcal B}(d)$ has rank $2r$ at
every distance. This is the rank of an operator on content coordinates;
the positional Gram instead measures shared dependence on distance.

\begin{proposition}[Content matching in a slow block]
\label{prop:slow-content-matching}
Let $H>0$ and $\eta=H\max_{k\in\mathcal B}|\omega_k|$. For $|d|\le H$,
\begin{equation}
\left|q_{\mathcal B}^{\top}\mathcal R_{\mathcal B}(d)k_{\mathcal B}
-q_{\mathcal B}^{\top}k_{\mathcal B}\right|
\le\eta\|q_{\mathcal B}\|_2\|k_{\mathcal B}\|_2.
\label{eq:slow-content-bound}
\end{equation}
For $2r$ content symbols represented by distinct coordinate unit vectors,
let the query and correct key share a vector and all incorrect keys use
other vectors. Their unscaled score margin is at least $1-2\eta$ across
arbitrary key distances in $[-H,H]$.
\end{proposition}

\begin{proof}
For each pair,
$\|R(\omega d)-I\|_2=2|\sin(\omega d/2)|\le|\omega d|$.
The block-diagonal norm is the maximum pair norm. The bilinear norm bound
gives Eq.~\eqref{eq:slow-content-bound}. Unrotated matching scores are one
for the correct symbol and zero for an incorrect symbol. Each score
changes by at most $\eta$, giving the margin bound. In particular,
$\eta<1/2$ preserves the correct ranking.
\end{proof}

As distinct frequencies approach zero, their positional subspaces approach
$\operatorname{span}\{1,d\}$ while the content construction remains
distinct in all $2r$ coordinates. This provides a concrete distinction
between repeated positional dependence and repeated content information,
consistent with prior analyses of frequency use \citep{barbero2025round}.
The construction concerns specified content vectors and the small-phase
regime; the next subsection evaluates the complete finite frequency tables.

\subsection{Task gains with lower full-pair effective rank}
\label{sec:allocation-rank-quality}
We use the Llama mathematical native grid
$\omega_k^N=500000^{-k/64}$, native length $8192$, $s=4$ and band $[18,35]$.
The TailSpline and MrPro formulas determine all frequencies without model
observations. Effective rank uses the full-pair Gram from
Eq.~\eqref{eq:cross-gram}, with uniform separations on $[0,H]$.

\begin{table}[ht]
\centering\small
\caption{\textbf{A successful allocation can have lower positional effective rank.}
Ranks describe the full mathematical frequency tables; task scores reuse
the clean paired evaluations in Table~\ref{tab:clean-length-main}.
The $16$K and $32$K task panels contain $650$ and $2600$ paired prompts.}
\label{tab:allocation-rank-quality}
\begin{tabular}{@{}lrrrr@{}}
\toprule
Window & $r_2$ TailSpline & $r_2$ MrPro & RULER TailSpline (\%) & RULER MrPro (\%)\\
\midrule
$16$K & $8.283391$ & $8.739887$ & $86.0949$ & $82.7051$\\
$32$K & $10.077270$ & $10.214888$ & $68.2660$ & $56.5436$\\
\bottomrule
\end{tabular}
\end{table}

Two independent continuous-Gram implementations agree within
$3.82\times10^{-12}$. Explicit integer-position features over
$0,\ldots,H-1$, an integer closed form, FP32 table rounding, and causal
pair weights proportional to $H-d$ preserve the rank ordering at
$8/16/32$K. The latter weights are a declared separation measure.
At $8$K the continuous ranks are $7.068167/7.716063$ for T/P.
These observations establish task gains with lower normalized positional
diversity; they do not identify rank reduction as the cause of the gains.

\subsection{How interior wavelengths respond to extension scale}
\label{sec:allocation-scale-response}
Write $\nu_k(s)=\omega_k^N s^{-m_k(s)}$ and let
$A_k(s)=\omega_k^N/\nu_k(s)$ be the wavelength multiplier.
For the index ramp in the released YaRN implementation
\citep{peng2024yarn}, with fixed interior coordinate $0<t<1$,
\begin{equation}
A_Y(s,t)=\frac{1}{1-t+t/s},\qquad
\lim_{s\to\infty}A_Y(s,t)=\frac{1}{1-t}.
\label{eq:yarn-interior-scale-limit}
\end{equation}
A fixed frequency blend $\gamma+(1-\gamma)/s$ similarly has limit
$1/\gamma$ for $\gamma>0$; the fully interpolated tail $\gamma=0$ grows
with $s$. MrPro and TailSpline instead have
$A_P=s^{P_q}$ and $A_T=s^{T_q}$ for their fixed finite-grid exponents.
Their log-scale derivatives are
\begin{equation}
\frac{\partial\log A_Y}{\partial\log s}
=\frac{t}{s(1-t)+t},\qquad
\frac{\partial\log A_P}{\partial\log s}=P_q,\qquad
\frac{\partial\log A_T}{\partial\log s}=T_q.
\label{eq:allocation-scale-derivatives}
\end{equation}
These identities follow by differentiation. They describe constructor
responses to $s$, rather than task-performance guarantees.

Index and turn-count ramps use different blend weights. All numerical
comparisons here use one shared band and the index-ramp equation of the
released YaRN implementation (revision \texttt{995db5b}).
They are frequency-table comparisons, not a re-evaluation of the YaRN
baseline runs in the MrRoPE paper.

\paragraph{Magnitude of the TailSpline intervention.}
Direct subtraction of the finite-grid formulas gives
\begin{equation}
T_q-P_q=\frac{q(n-q)(3n+q+1)}{n(n+1)(2n+1)}\ge0,
\qquad \frac{A_T}{A_P}=s^{T_q-P_q}.
\label{eq:tailspline-wavelength-order}
\end{equation}
For the Llama $n=17$ grid, the maximum wavelength ratio is
$1.328763$, $1.765612$, $2.346080$ and $3.117385$ at $s=2,4,8,16$.
At pair $k=26$ ($q=8$), the multipliers are:
\begin{center}
\begin{tabular}{@{}lrrr@{}}
\toprule
Scale & YaRN index blend & MrPro & TailSpline\\
\midrule
$4$ & $1.545455$ & $1.385674$ & $2.423868$\\
$16$ & $1.789474$ & $1.920093$ & $5.875134$\\
\bottomrule
\end{tabular}
\end{center}
The YaRN/MrPro ordering at an interior pair can therefore change with
scale. TailSpline redistributes substantial wavelength growth throughout
the transition, in addition to its smaller terminal extra-gap jump.

\subsection{Two reference responses and finite phase intervals}
\label{sec:allocation-reference-intervals}
The native reference is $\omega$ and the fully interpolated reference is
$\omega/s$. For a declared phase tolerance $0<\theta\le\pi$, define
\begin{equation}
D_N(\nu;\theta)=\frac{\theta}{|\omega-\nu|},\qquad
D_D(\nu;\theta)=\frac{\theta}{|\nu-\omega/s|},
\label{eq:allocation-reference-distances}
\end{equation}
with infinity at a zero denominator. These are the maximal intervals
starting at zero on which the respective unwrapped phase deviation is at
most $\theta$. Since $\omega/s\le\nu_T\le\nu_P\le\omega$,
\begin{equation}
D_D(\nu_T;\theta)\ge D_D(\nu_P;\theta),\qquad
D_N(\nu_T;\theta)\le D_N(\nu_P;\theta).
\label{eq:allocation-reference-order}
\end{equation}
The exact operator difference is
$2|\sin(d(\nu-\nu_{\rm ref})/2)|$.
Thus the distances quantify two specified phase references; they are not
effective context lengths of a model.

For a fixed content component
$f_\nu(d)=a\cos(\nu d-\omega d_0)$, with $a>0$ and desired stretched
distance $d=sd_0$, its residual phase is $\delta=d(\nu-\omega/s)$.
If $0\le\delta_T\le\delta_P\le\pi$, then $f_T(d)\ge f_P(d)$ by the
monotonicity of cosine on $[0,\pi]$. The content phase specifies the
reference calculation. Outside that interval, reducing residual phase
from $2\pi$ to $\pi$ instead changes the response from $a$ to $-a$.
The conditional example explains a possible use of distance scaling,
while the full-model comparisons measure its actual task value.

\paragraph{Equal displacement does not imply a small phase intervention.}
On the same Llama $s=4$ mathematical grid, the largest T/C unwrapped
phase difference over integer distances $0,\ldots,32767$ is
$10.672514$ radians. The corresponding T/P value is $108.197815$.
For finite comparisons one can use
\begin{equation}
\rho_H(\delta)=\frac4H\sum_{d=0}^{H-1}\sin^2(d\delta/2)
=2-2\frac{\sin(H\delta/2)}{H\sin(\delta/2)}
\cos((H-1)\delta/2),
\label{eq:allocation-finite-phase-cost}
\end{equation}
with removable singularities evaluated by continuity. The fixed-slot
rotation RMS is $[K^{-1}\sum_k\rho_H(\nu_k-\widetilde\nu_k)]^{1/2}$.
At $H=32768$, it is $0.448329$ for T/C and $0.687543$ for T/P.
These values compare unscaled rotary kernels, before the shared
cosine/sine gain and the attention normalization.
For $H|\delta|\ll1$,
$\rho_H(\delta)=((H-1)(2H-1)/6)\delta^2+O(H^4\delta^4)$.
The small parameter is therefore distance times frequency difference;
matching total displacement alone does not supply it.
